\section{Introduction}
An alarm in a robotic system must identify a deviation and remain active for as long as the monitored condition warrants it. These decisions need different evidence. A predictor that follows recent measurements may produce a small error while the robot is still displaced from its approved trajectory. A fixed reference can expose that displacement, but errors in its plant model can also cause alarms during legitimate motion.

We study whether a separate return criterion can preserve alarm evidence through recovery without increasing alarms during authorized command changes. A following observer supplies a prediction error for alarm entry. A reference driven by the approved command supplies the return condition. Once triggered, the alarm clears only after the measured state stays close to this reference for ten consecutive samples.

We evaluate this rule on 144 simulated episodes and compare 23 methods and ablations. All context-aware methods receive the same protected commands and noisy joint observations. Two comparisons separate the effects of calibration and recovery logic: one calibrates each complete decision rule on the same normal data, while the other holds the entry threshold fixed. The experiment tests an engineering combination of established components; it makes no claim of algorithmic originality.

\section{Related work}
CASPER detects motion anomalies using external inertial measurements from a robotic arm and evaluates deviations introduced through joint-velocity changes \cite{casper}. Its successor investigates real-time detection on edge hardware \cite{edge}. Both concern movement monitoring, but their external-IMU inputs differ from the four joint-state channels and protected commands used here. Their published results therefore cannot be compared directly with this simulation.

TWINE combines a manufacturing digital twin with explainable anomaly detection \cite{twine}. Its supervised classification workflow uses labels that are unavailable to our normal-only procedure. Ayankoso et al. combine trajectory-change recognition with condition monitoring \cite{ayankoso}. Recognizing a changed trajectory is relevant to legitimate operation, although the movement alone does not establish whether a command was authorized.

Observer-based fault diagnosis already addresses the effects of gain and model uncertainty on detection \cite{meseguer}. CUSUM accumulates evidence of persistent changes \cite{nist}. More recent work examines adaptive interval observers for actuator faults \cite{rapio} and digital-twin-based false-data-injection detection in robotic vehicles \cite{kaisar}. Full-method comparisons with these two recent studies remain unresolved because complete manuscripts or implementation details were unavailable in the source audit.

Our executable comparisons use specified predictors and classical alarm rules under common information access. The earlier nine-channel IMU predictors remain a separate experiment. No published detector's complete experimental protocol has been reproduced in the present study.
